% from ega2-errata-addenda-fr.tex
\oldpage[II]{221}
\markboth{ERRATA ET ADDENDA, LISTE 1}{ERRATA ET ADDENDA, LISTE 1}
annelés en anneaux locaux $u:X\to Y$ et d'un $u$-morphisme
$h:\sh{G}\to\sh{F}$ de Modules; ces
morphismes (pour $(X,\sh{F})$ et $(Y,\sh{G})$ fixés) forment un ensemble
que nous noterons $\Hom((X,\sh{F}),(Y,\sh{G}))$; l'application
$(u,h)\mapsto(\rho'(u),\Gamma(h))$ est une application canonique
\begin{equation}
\label{ega2:addendum:I.1.8.9.1-fr}
\begin{split}
  \Hom((X,\sh{F}),(Y,\sh{G}))\to
  \Hom\bigl(
    &(\Gamma(Y,\sh{O}_Y),\Gamma(Y,\sh{G})),\\
    &(\Gamma(X,\sh{O}_X),\Gamma(X,\sh{F}))
  \bigr),
\end{split}
\tag{1.8.9.1}
\end{equation}
fonctorielle en $(X,\sh{F})$ et $(Y,\sh{G})$, le second membre étant
l'ensemble des di-homomorphismes correspondant aux anneaux et aux
modules considérés (0, 1.0.2).
\end{env}

\begin{corollary}[1.8.10]
\label{ega2:addendum:I.1.8.10-fr}
Soient $Y$ un espace annelé en anneaux locaux, $\sh{G}$ un
$\sh{O}_Y$-Module. Pour que $Y$ soit un schéma affine et $\sh{G}$ un
$\sh{O}_Y$-Module quasi-cohérent, il faut et il suffit que pour tout
couple $(X,\sh{F})$ formé d'un espace annelé en anneaux locaux $X$ et
d'un $\sh{O}_X$-Module $\sh{F}$, l'application canonique (1.8.8.1) soit
bijective.
\end{corollary}

Nous laissons au lecteur le soin de développer le raisonnement, calqué
sur celui de (1.8.3), et où on utilise (1.8.1) et (1.8.7).

\begin{remark}[1.8.11]
\label{ega2:addendum:I.1.8.11-fr}
Les énoncés (I, 1.7.3), (I, 1.7.4) et (I, 2.2.4) sont des cas
particuliers de (1.8.1), ainsi que la définition (I, 1.6.1); de même
(I, 2.2.5) résulte de (1.8.7). La proposition (1.8.7) entraîne aussi
(I, 1.6.3) (et par suite (I, 1.6.4)) comme cas particulier, car si $X$
est affine et $\Gamma(X,\sh{F})=N$, les foncteurs
\[
  M\longmapsto
  \Hom_{\sh{O}_S}(\widetilde M,u_*(\widetilde N))
  \quad\hbox{et}\quad
  M\longmapsto
  \Hom_{\sh{O}_S}(\widetilde M,(N_{[\vphi]})\supertilde)
\]
(où $\vphi:A\to\Gamma(X,\sh{O}_X)$ correspond à $u$) sont isomorphes
par (1.8.7) et (I, 1.3.8). Enfin, (I, 1.6.5) (et par suite (I, 1.6.6))
résulte de (1.8.6) et du fait que pour tout $f\in A$, les $A_f$-modules
\[
  N'\otimes_{A'}A_f
  \quad\hbox{et}\quad
  (N'\otimes_{A'}A)_f
\]
(notations de (I, 1.6.5)) sont canoniquement isomorphes.
\end{remark}

\paragraph{Insérer après (I, 3.2.8) :}

\begin{env}[3.2.9]
\label{ega2:addendum:I.3.2.9-fr}
Il résulte de (1.8.1) que l'on peut préciser (I, 3.2.2) de la façon
suivante : $Z=\Spec(B\otimes_A C)$ est non seulement un produit de
$X=\Spec(B)$ et de $Y=\Spec(C)$ dans la catégorie des $S$-préschémas,
mais aussi dans la catégorie des espaces annelés en anneaux locaux,
au-dessus de $S$ (avec une définition des $S$-morphismes calquée sur
celle de (I, 2.5.2)). La démonstration de (I, 3.2.6) prouve alors en
fait que pour deux $S$-préschémas quelconques $X,Y$, le préschéma
$X\times_S Y$ est non seulement le produit de $X$ et $Y$ dans la
catégorie des $S$-préschémas, mais encore dans la catégorie des espaces
annelés en anneaux locaux au-dessus du préschéma $S$.
\end{env}

\paragraph{(I, 4.4.5).}
Ligne 14 du bas de la p.~126, remplacer $B$ par $A$; ligne 13 du bas,
remplacer «~$A$-algèbre~» par «~$B$-algèbre~».

\paragraph{(I, 5.3.5).}
Ligne 2 du bas de la p.~132, remplacer $Y\to S$ par $g:Y\to S$.

\paragraph{(I, 6.3.2.1).}
Ligne 6 du bas de la p.~144, lire :
\[
  D(g_i)\subset W.
\]

\paragraph{(I, 7.3.8).}
Remplacer le texte actuel, qui est erroné, par le suivant :

Soient $X,Y$ deux préschémas intègres, ce qui entraîne que
$\sh{R}(X)$ (resp. $\sh{R}(Y)$) est un $\sh{O}_X$-Module (resp. un
$\sh{O}_Y$-Module) quasi-cohérent (I, 7.3.3). Soit $f:X\to Y$ un
morphisme dominant; alors il existe un homomorphisme canonique de
$\sh{O}_X$-Modules
